\documentclass[11pt,a4paper]{article}
\usepackage[T1]{fontenc}
\usepackage{amsmath,amssymb,amsthm}
\usepackage{needspace}
\usepackage[margin=27mm]{geometry}
\usepackage[hidelinks]{hyperref}
\hypersetup{pdftitle={Interaction support in smooth-number orders and corner-cut fans},
 pdfauthor={},pdfsubject={Anonymous working draft for mathematical review}}
\newtheorem{theorem}{Theorem}
\newtheorem{lemma}[theorem]{Lemma}
\newtheorem{corollary}[theorem]{Corollary}
\newtheorem{proposition}[theorem]{Proposition}
\theoremstyle{remark}\newtheorem{remark}[theorem]{Remark}
\newcommand{\NN}{\mathbb Z_{\geq0}}
\newcommand{\RR}{\mathbb R}
\newcommand{\ZZ}{\mathbb Z}
\newcommand{\supp}{\operatorname{supp}}
\newcommand{\card}[1]{\lvert#1\rvert}
\newcommand{\fl}[1]{\lfloor#1\rfloor}
\newcommand{\ce}[1]{\lceil#1\rceil}
\newcommand{\ga}{\gamma}
\title{Interaction support in smooth-number orders\\and corner-cut fans}
\author{}
\date{19 October 2026}
\begin{document}
\maketitle
\begin{center}\small Anonymous working draft for mathematical review.\end{center}
\begin{abstract}
Order the nonnegative integer lattice by a positive weight vector.
How many coordinates must a comparison involve before it can first
affect a prefix of length \(n\)? We prove that the minimum first tied
rank of a generic comparison wall with support \(s\) is
\(\Theta((3\sqrt3/2)^s/\sqrt{s})\), with absolute constants.
Consequently the worst-case coordinate-face dimension needed to
reconstruct a labeled \(n\)-term ordered prefix is
\[
 \frac{\log n+\tfrac12\log\log n}{\log(3\sqrt3/2)}+O(1).
\]
The same parameter is the maximum support of a positive-weight facet
of a prefix cone and of a sufficiently generic \(n\)-point ideal's
state fan, uniformly over ambient dimension. It yields sparse linear
certificates for a realizable prefix. Finite prime approximation
transfers the face-information result to smooth-number exponent
orders. The general Hilbert-zonotope framework has a different
support scale: arbitrary length-\(n\) quotient algebras admit walls
with \(2\log_2 n+O(1)\) coordinates. We give an explicit boundary
example and separate the structural statements from running-time
claims.
\end{abstract}

\section{The question and its relation to exponent streams}
Let \(d,n\geq1\), and let \(w\in\RR_{>0}^d\) have rationally independent coordinates.
Write \(W_n(w)\) for the first \(n\) vectors of \(\NN^d\), in increasing
order of \(w\cdot x\), with the zero vector at rank \(1\).
For distinct primes \(p_1,\ldots,p_d\), taking \(w_i=\log p_i\)
gives the exponent vectors of the fixed-basis smooth numbers.
Coordinates retain their labels; numerical weight order is not
additional input. In the prime model the labels may have any fixed
order, and the numerical prime values are not supplied as additional
data in the reconstruction problem.

For a nonempty coordinate set \(S\), the \emph{face word}
\(W_n(w_S)\) orders the lattice on \(S\). It contains \(n\) vectors
from that face, not the first \(n\) entries of a coordinate projection
of the full word. Let \(R_{\rm all}(n)\) be the least integer \(r\)
such that all labeled face words of length \(n\) on at most \(r\)
coordinates determine \(W_n(w)\), uniformly in \(d\) and \(w\).
Let \(R_{\rm prime}(n)\) restrict the realizing weights to distinct
prime logarithms. Define both parameters to be \(0\) at \(n=1\).

This question comes from the deletion and occurrence structure of
the individual exponent streams. If \(E_i(j)\) is coordinate \(i\)
of the \(j\)-th vector, define \(T(E_i)(j)\) as the number of earlier
occurrences of the same symbol. Fixing coordinate \(i\) leaves the
order on the other coordinates. Therefore
\[
 T(E_i)(j)=\rho_i(x_j\text{ with coordinate }i\text{ deleted})-1,
 \qquad T^2(E_i)=E_i,
\]
where \(\rho_i\) is the one-based rank on the complementary face.
Indeed occurrences of each face rank have clock exponents
\(0,1,2,\ldots\). Deleting zero entries of \(E_i\) and subtracting
one corresponds to translating lattice points by the \(i\)-th unit
vector, and reproduces the full order. Placement arrays and
signature-sequence duality are established background
\cite{KB04,K22}. For two generators the partner is the other
exponent; in higher dimension it encodes the rank of an entire
face vector. Our concern is the finite information needed to
recover the ordering hidden in these ranks.

Put \(\ga=3\sqrt3/2\). All unqualified logarithms below are natural.

\begin{theorem}[Sharp face information]\label{thm:faces}
As \(n\to\infty\),
\[
 R_{\rm all}(n)=R_{\rm prime}(n)
 =\frac{\log n+\tfrac12\log\log n}{\log\ga}+O(1).
\]
The equality of the two parameters is exact for each \(n\);
the displayed asymptotic determines their value to bounded
additive error. Constants are independent of ambient dimension
and of the prime bases.
\end{theorem}

The statements concern ordinal exponent data. An arithmetic
algorithm supplied with the primes can compare numerical products
directly. Nor is requesting all faces cheap: their number may be
\(\sum_{j\le r}\binom dj\), each containing \(n\) vectors.
No random-access running-time improvement follows from
Theorem~\ref{thm:faces} alone.

\section{Comparison walls and a composition estimate}
For a primitive mixed-sign \(z\in\ZZ^d\), put \(u=z_+\) and
\(v=z_-\). A \emph{generic positive wall point} is \(w>0\) such that
\[
 \{a\in\ZZ^d:w\cdot a=0\}=\ZZ z.
\]
At height \(h=w\cdot u=w\cdot v\), the only tied points are
\(u,v\). To see this, a point at that height differs from \(u\)
by an integer multiple of \(z\); nonnegativity permits only the
two endpoints. The lowest tie anywhere in the lattice is this pair:
every other tied pair is a nonnegative translate of positive integer
multiples of \(u,v\).
Its \emph{first tied rank} is \(1+\card{\{x:w\cdot x<h\}}\).

Let \(J(s)\) be the minimum first tied rank of a generic positive
wall whose normal has support \(s\). One may work in dimension \(s\):
additional coordinates can only increase the count, and can be
chosen large when proving an upper bound.

\begin{lemma}[Box bound]\label{lem:box}
At a generic positive wall,
\[
 J\geq \frac12\prod_i(1+|z_i|)\geq2^{s-1}.
\]
\end{lemma}
\begin{proof}
Pair the points \(x\) and \(u+v-x\) in the integer box
\(0\leq x\leq u+v\). Their scores add to \(2h\), and only \(u,v\)
have score \(h\). Thus exactly
\((\prod_i(1+|z_i|)-2)/2\) box points lie strictly below \(h\).
The product is even because a primitive vector has an odd
coordinate.
\end{proof}

The box and flip viewpoints are classical ingredients
\cite{M99,CCS}. The additional points with repeated coordinates
are essential for the sharper bound.

\begin{lemma}[Strict composition halfspaces]\label{lem:moment}
Let \(s\geq2\), \(q=\fl{s/2}\), and let \(X\) be uniform on
\(\{x\in\NN^s:\sum_i x_i=q\}\). If \(\beta\ne0\) and
\(\sum_i\beta_i=0\), then, for \(Y=\beta\cdot X\),
\[
 \mathbb E Y^4\leq12(\mathbb E Y^2)^2,\qquad
 \Pr(Y<0),\Pr(Y>0)\geq1/48.
\]
\end{lemma}
\begin{proof}
Use falling factorials \((q)_j\), rising factorials \(s^{(j)}\),
and \(u_j=(q)_j/s^{(j)}\), interpreted as zero for \(j>q\).
The standard factorial-moment identity is
\[
 \mathbb E\prod_i(X_i)_{a_i}
 =\frac{(q)_{\sum a_i}\prod_i a_i!}{s^{(\sum a_i)}}.
\]
For a self-contained derivation, the numerator's generating function is
\[
 \prod_i(a_i!)t^{\sum a_i}(1-t)^{-s-\sum a_i}.
\]
Take its coefficient of \(t^q\) and divide by
\(\binom{s+q-1}{q}\). This is also the all-ones-parameter
Dirichlet--multinomial specialization of the published general
formula \cite[Section 4.2 and Appendix A]{T14}.

Put \(A=\sum_i\beta_i^2\) and \(C=\sum_i\beta_i^4\leq A^2\).
Expanding powers with the factorial moments and
\(\sum_i\beta_i=0\) gives
\begin{align*}
 \mathbb EY^2&=(u_1+u_2)A,\\
 \mathbb EY^4&=(u_1+7u_2+12u_3+6u_4)C
                +(3u_2+6u_3+3u_4)A^2.
\end{align*}
Set \(r=q/s\) and \(t=u_2\). Then
\(1/3\leq r\leq1/2\), \(0\leq t\leq r^2\),
\(u_3\leq rt\), and \(u_4\leq r^2t\). Hence
\[
 \frac{\mathbb EY^4}{A^2}
 \leq r+(10+18r+9r^2)t\leq12(r+t)^2.
\]
For the last inequality, the difference between right and left
is decreasing in \(t\) on the indicated interval: its derivative
is at most \(6r+15r^2-10<0\). At \(t=r^2\) the difference is
\(r(2r+6r^2+3r^3-1)>0\).

Write \(\sigma^2=\mathbb EY^2>0\). H\"older interpolation gives
\(\mathbb E|Y|\geq\sigma^3/\sqrt{\mathbb EY^4}
\geq\sigma/\sqrt{12}\).
Since \(\mathbb EY=0\), \(\mathbb EY_-=\mathbb E|Y|/2\).
Cauchy--Schwarz bounds this by
\(\sigma\sqrt{\Pr(Y<0)}\), giving \(1/48\).
Apply the same argument to \(-Y\).
\end{proof}

No separate novelty claim is made for the factorial moments,
the fourth-moment method, or the particular nonoptimal constant.

\begin{lemma}[Strict lattice count]\label{lem:count}
If \(w_i>0\), \(h>0\), and \(\sum_{i=1}^s w_i\leq2h\), then
\[
 \card{\{x\in\NN^s:w\cdot x<h\}}
 \geq \frac1{48}\binom{s+\fl{s/2}-1}{\fl{s/2}}.
\]
\end{lemma}
\begin{proof}
Put \(q=\fl{s/2}\). If the weights are not all equal, apply
Lemma~\ref{lem:moment} to
\(\beta_i=w_i-\sum_jw_j/s\). At least the stated number
of degree-\(q\) points are strictly below their mean
\(q\sum_iw_i/s\leq h\).
If all weights equal \(a\), the entire layer is strictly
below \(h\) unless \(qa=h\). In that exception \(s=2q\);
the degree-\(q-1\) layer is below \(h\) and has
\[
 \binom{s+q-2}{q-1}
 =\binom{s+q-1}{q}\frac q{s+q-1}
 \geq\frac13\binom{s+q-1}{q}.
\]
This handles the zero-variance case separately.
\end{proof}

\section{The sharp interaction scale}
\begin{theorem}\label{thm:J}
Uniformly for \(s\geq2\),
\[
 J(s)=\Theta(\ga^s/\sqrt{s}).
\]
For every prescribed finite \(M\), the upper constructions can be
perturbed to opposite sides of a wall while preserving all proper
coordinate-face words through \(M\) terms.
\end{theorem}
\begin{proof}
At a wall with support \(s\), each nonzero \(|z_i|\geq1\), so
\(\sum_{i\in\supp z}w_i\leq w\cdot(u+v)=2h\).
Lemma~\ref{lem:count} and Stirling's formula give the lower bound,
since
\[
 \log\binom{s+\fl{s/2}-1}{\fl{s/2}}
   =s\log\ga-\tfrac12\log s+O(1).
\]

For the upper bound put \(a=\fl{s/2}\), \(b=\ce{s/2}\),
and \(z=(1^a,-1^b)\). Start with weights \(b/a\) in the
first group and \(1\) in the second. The common height is \(b\).
Choose a generic point on this wall within
\(\eta<1/(s+1)\) of these weights. Its common height is at most
\(b+a\eta\), whereas every point of degree at least \(b+1\)
has score at least
\((b+1)(1-\eta)>b+a\eta\).
Therefore
\[
 J(s)\leq\binom{s+b}{b}=\Theta(\ga^s/\sqrt{s}).
\]
Such generic points exist by avoiding the countably many
other rational hyperplanes inside a relatively open wall patch.

Every proper coordinate subvector at that point is independent:
a relation on a proper face cannot be a nonzero multiple of
the full-support normal \(z\). Each prescribed finite face word
is determined by finitely many strict inequalities. A sufficiently
small perturbation of the last weight in either direction preserves
all of them. It reverses the lowest full tie, changing the full word
at its first tied rank. Perturbations can be chosen off all rational
hyperplanes.
\end{proof}

\begin{lemma}[Finite prime realization]\label{lem:prime}
Any finite collection of strict positive-weight comparisons can
be preserved by distinct prime logarithms. Two perturbations
differing only in the last coordinate can share all their other
prime bases.
\end{lemma}
\begin{proof}
For large \(L\), Bertrand's postulate supplies a prime in
\((e^{Lw_i},2e^{Lw_i})\). The normalized logarithmic error is
at most \(\log2/L\). All finitely many strict comparisons
persist for large \(L\). Choose the target vectors with distinct
coordinates; their intervals become disjoint, giving distinct
bases. Shared coordinates use the same intervals and primes.
\end{proof}
Finite prime realization is established in
\cite[Proposition 6.2]{V02}; the lemma records the elementary
version used here, including shared coordinates for the two
perturbations. It gives no new prime-distribution estimate.

\section{Exact inversion and reconstruction}
A finite prefix's weight region is an open polyhedral cone.
Indeed impose its consecutive comparisons, positivity, and the
comparison of its last point with every coordinatewise minimal
excluded point. The excluded border is finite, and every excluded
point dominates a border point.

\begin{theorem}[Exact interaction parameter]\label{thm:kappa}
For \(n\geq2\), let
\(\kappa(n)=\max\{s\geq2:J(s)\leq n\}\).
Set \(\kappa(1)=0\).
Then
\[
 R_{\rm all}(n)=R_{\rm prime}(n)=\kappa(n).
\]
\end{theorem}
\begin{proof}
Take two weights giving the same supplied face words on at most
\(\kappa(n)\) coordinates. The intersection of the corresponding
face-prefix cones is open and convex. Join the weights by a
generic path inside it. A finite arrangement suffices: every point
that can appear or first tie by rank \(n\) has every coordinate
at most \(n\), since larger values give too many smaller pure powers.
The path can avoid the countably many intersections of distinct
rational comparison hyperplanes. The finite relevant arrangement
then detects every prefix change.

If the full prefix changes, a wall with primitive normal \(z\)
has its lowest tied rank at most \(n\). Restricting to its support
\(S\) cannot increase this rank, so \(J(|S|)\leq n\).
Its face is among the supplied data, and crossing the wall changes
that face's length-\(n\) word. This contradicts the common cone.

Conversely, a generic \(s\)-coordinate wall attaining \(J(s)\leq n\)
has perturbations agreeing through \(M=n\) on all proper faces and
differing in the full prefix. Lemma~\ref{lem:prime} transfers the
two finite patterns to prime bases. Thus fewer than \(s\) coordinates
cannot suffice even in the prime class. The box bound makes
\(\kappa(n)\) finite.
\end{proof}

\begin{proof}[Proof of Theorem~\ref{thm:faces}]
The lower estimate for \(J(s)\) implies that any \(J(s)\leq n\)
satisfies
\[
 s\log\ga-\tfrac12\log s\leq\log n+O(1).
\]
The box bound first gives \(s=O(\log n)\), yielding the stated
upper bound. For the reverse inequality use the even-dimensional
upper construction with \(s=2k\) and first rank at most
\(\binom{3k}{k}\). Stirling gives
\(\log\binom{3k}{k}=2k\log\ga-\tfrac12\log k+O(1)\).
Taking
\[
 k=\fl{\frac{\log n+\tfrac12\log\log n}{2\log\ga}-C}
\]
for a sufficiently large fixed \(C\) proves the matching lower bound.
\end{proof}

\section{Sparse validity cones and interval certificates}
\begin{theorem}\label{thm:cone}
For \(n\geq2\), every positive-weight facet of a realizable length-\(n\) prefix
cone has support at most \(\kappa(n)\). Some cone needs this support.
The same bound is sufficient for a finite linear inequality
description, retaining coordinate positivity separately.
\end{theorem}
\begin{proof}
Choose a generic positive relative interior point of a facet.
Crossing it changes the prefix, so its normal \(z\) has
\(J(|\supp z|)\leq n\), as above. Coordinate boundary facets
are represented by positivity. In a finite halfspace description
each facet has a defining normal. Retaining these normals defines
the same closed cone. Every other valid normal is a nonzero
nonnegative combination of them, so their strict inequalities
also imply all original strict inequalities. Realizability
provides the strict interior point needed here.

Conversely a wall attaining \(J(s)\leq n\) supplies a genuine
prefix facet of support \(s\). Any finite inequality description
must contain a normal parallel to that facet; normals with
smaller support cannot define it.
\end{proof}

An explicit sufficient cutoff follows from the proof, without
computing \(J(s)\). Put \(p=\fl{(s-1)/2}\), \(q=\fl{s/2}\), and
\[
 B_s=\max\left(2^{s-1},\
  1+\ce{\binom{s+p}{p}/s},\
  1+\ce{\binom{s+q-1}{q}/48}\right).
\]
For the middle bound, each cyclic permutation orbit of the
points of degree at most \(p\) has average height strictly
below \(h\), so supplies one point below \(h\).
Thus \(J(s)\geq B_s\), and
\(r(n)=\max\{s:B_s\leq n\}\) is sufficient.
One may improve this with the elementary exact values
\(J(2)=2,J(3)=5,J(4)=10\), proved in the research artifact.

For an actual certificate, first generate a proposed word at
exact positive rational weights and replay through \(n+1\)
points to verify a strict prefix and boundary. Form the usual
consecutive and minimal-border comparisons. Delete all
nonpositivity normals above a proved sufficient support cutoff.
For a rational box \(w_i\in[l_i,u_i]\), \(l_i>0\), check
\[
 \sum_{z_i\geq0}z_i l_i+\sum_{z_i<0}z_i u_i>0
\]
for every retained normal \(z\). Theorem~\ref{thm:cone}
certifies the whole box, including any prime logarithms
rigorously enclosed in it. It is essential to establish a strict
realization first; deletion need not preserve the feasibility
question for an arbitrary impossible word. This procedure reads
the whole prefix and border. It makes no running-time claim
against the published recognition algorithms.

\section{Generic state fans and the quotient boundary}
Work over \(\mathbb C\) in the algebraic statements; normal cones
use minimization. A state-fan
facet means a codimension-one wall between full-dimensional
positive-weight regions. For a length-\(n\) ideal, the state polyhedron and its normal fan
encode its initial ideals. The published framework
\cite{OS99,BOT03} identifies the state polyhedron of a sufficiently
generic \(n\)-point configuration with the corner-cut polyhedron
\[
 P_n^d=\operatorname{conv}\left\{\sum_{x\in D}x:
                     D\subset\NN^d,\ |D|=n\right\}.
\]
A precise sufficient genericity condition is that every
\(n\)-element subset of
\(V_n^d=\{v:\prod_i(v_i+1)\leq n\}\) is a quotient basis.
This identification is prior art.

\begin{corollary}\label{cor:generic}
For \(n\geq2\), the maximum support of a state-fan facet meeting the positive
weight cone, over all ambient dimensions and sufficiently
generic \(n\)-point configurations, equals \(\kappa(n)\).
\end{corollary}
\begin{proof}
A corner-cut cone is refined by the cones for its internal
ordered prefixes. A positive-weight wall therefore has a
primitive normal whose lowest tied rank is at most \(n\);
its support is at most \(\kappa(n)\).

For the reverse bound, take a wall attaining \(J(s)\leq n\)
at height \(h\). Append \(n-J(s)\) new coordinate weights in
\((\max(h/2,h-\min_iw_i),h)\). Each contributes exactly its
unit vector below \(h\). Squares and sums with old nonzero
points or other new coordinates are above \(h\).
Choose the added weights generically, preserving only the
original relation. The first tied rank is now precisely \(n\).
Across this wall the two \(n\)-point cut sums differ by \(v-u\),
giving a positive-weight facet with support \(s\).
The published generic-configuration identification transfers it
to the state fan.
\end{proof}
For \(n=1\) the prefix is the zero vector and there is no
positive-weight facet in either setting; no maximum over an
empty family of facets is required.

The corresponding statement for arbitrary quotient algebras has
a different scale. Let \(A(n)\) be the maximum support of a
noncoordinate state-fan facet meeting positive weights, over
all length-\(n\) ideals and ambient dimensions.

\begin{proposition}[Boundary for arbitrary quotients]\label{prop:quotient}
For \(n\geq2\),
\[
 2\fl{\log_2(n+2)}-2\leq A(n)\leq2\fl{\log_2 n}.
\]
\end{proposition}
\begin{proof}
The upper bound is a direct consequence of the
Hilbert-zonotope refinement \cite[Proposition 3.2 and
Theorem 3.5]{BOT03}. Every refining direction is a difference
of two members of \(V_n^d\), each with support at most
\(\fl{\log_2 n}\). A facet of the coarser state fan must
lie in one of the refining wall hyperplanes.

For the lower bound put \(k=\fl{\log_2(n+2)}-1\),
\(U=x_1\cdots x_k\), \(V=y_1\cdots y_k\), and
\(\ell=n-(2^{k+1}-2)\). Let
\[
 S=\operatorname{div}(U)\cup\operatorname{div}(V)
            \cup\{t,t^2,\ldots,t^\ell\}.
\]
Let \(K\) be the monomial ideal of the complement of \(S\)
and set \(I=K+(U-V)\). The set \(S\) has \(n+1\) members.
Every variable annihilates \(U,V\) modulo \(K\);
hence their relation removes exactly one dimension and
the quotient has length \(n\).

For either leading choice of \(U-V\), the monomial
generators of \(K\) and this binomial form a Gr\"obner basis.
Indeed an S-pair with a monomial generator reduces to a
nonconstant multiple of the other socle monomial, in \(K\).
Monomial--monomial S-pairs vanish. Thus the two initial
ideals are \(K+(U)\) and \(K+(V)\), separated at positive
weights by
\(\sum_iw_{x_i}=\sum_iw_{y_i}\), of support \(2k\).
\end{proof}
The general upper bound is credited as an immediate published
consequence. The example delimits Corollary~\ref{cor:generic};
it is not proposed as an independent universal-basis algorithm.
A length-\(6\) quotient already has a support-\(4\) wall,
whereas the full lattice cannot first have one before rank \(10\).

\section{Prior-art boundary, verification, and assessment}
Corner cuts, generic state polyhedra, and recognition via finite
linear inequalities are established \cite{OS99,BOT03}.
The recognition theorem of \cite[Section 3]{OS99} allows growing
dimension and uses Borel-fixedness to obtain a complete
frontier description. It assumes that the staircase itself is
given. Our face input does not give that staircase, and our
support bound is not a new recognition-time theorem.
Wagner's fixed-dimension counting and flag bounds \cite{W02}
address a different parameter. The finite-chain reductions of
\cite[Sections 5--6]{VP09} likewise use a known full chain and
its front; they can retain coprime comparisons of full support.
Boolean flips and minimal determining comparisons
\cite{M99,CCS} are credited; repeated powers are indispensable
to the lattice count here.

The edge conditions of \cite[Section 5]{Mu03} constrain the affine
span of exchanged cut points. Our generic wall exchanges only
two points, so that condition alone supplies no support/rank
count. The weighted-average degree bounds and greedy model
selection of \cite[Theorems 2--7]{B08} use the same established
geometry but a different objective. The lattice-point
decomposition question in \cite[Theorem 4, Corollary 5 and
Section 4]{L21} is also distinct from facet-normal support.
These are theorem-level comparisons with the inspected results,
not claims that no other derivation exists.

The candidate contribution is the sharp support/rank law,
its exact face/facet interpretation, and the generic-state
consequence, rather than those frameworks or the moment method.
The repository's bounded primary comparison records the
statements inspected and remaining priority limits. This
draft does not assert independent confirmation of novelty,
journal acceptance, or a new prime-distribution result.

The dependency-free research checker is
\texttt{scripts/run\_face\_reconstruction\_checks.py};
its report is
\texttt{results/research/face\_reconstruction\_checks.json}.
It checks exact moments on complete composition layers,
strict boundary cases, cyclic orbits, even and odd support
obstructions, prime witnesses, and sparse-cone certificates.
Separate exact product and pointer generators validate prime
prefixes. The quotient examples check every nontrivial
Buchberger pair for both leading choices. Large finite
dimension brackets use exact binomial integers. These checks
support the delicate steps; the uniform theorems depend
on the proofs, not a fitted asymptotic curve.

This is a structural result for growing interaction dimension.
Its practical algorithmic value remains to be established.
The single-stream rigidity and conditional finite-mask results
are preserved separately in the research note and are not
needed for the main theorem.

\Needspace{7\baselineskip}
\paragraph{Reproducibility.}
From the repository root run
\begin{quote}\ttfamily\small
python3 -B scripts/run\_face\_reconstruction\_checks.py\\
./build.sh\\
python3 tests/test\_smoke.py
\end{quote}
The core CLI, its independent rank auditor and previous
manuscripts are unchanged.

\paragraph{Assistance and review status.}
This anonymous working draft was developed with substantial
assistance from OpenAI's Codex in proof development,
computation and literature comparison. The proof package
has undergone internal checking; no independent referee
assessment is claimed. Authorship and submission decisions
have not been made.

\begin{thebibliography}{10}
\bibitem{OS99} S. Onn and B. Sturmfels,
\emph{Cutting corners}, Advances in Applied Mathematics 23 (1999),
29--48. \href{https://doi.org/10.1006/aama.1999.0645}{doi:10.1006/aama.1999.0645}.
\bibitem{BOT03} E. Babson, S. Onn and R. Thomas,
\emph{The Hilbert zonotope and a polynomial time algorithm for
universal Gr\"obner bases}, Advances in Applied Mathematics 30 (2003),
529--544. Inspected preprint:
\href{https://arxiv.org/abs/math/0207135}{arXiv:math/0207135}.
\bibitem{W02} U. Wagner, \emph{On the number of corner cuts},
Advances in Applied Mathematics 29 (2002), 152--161.
\href{https://doi.org/10.1016/S0196-8858(02)00014-3}{doi:10.1016/S0196-8858(02)00014-3}.
\bibitem{VP09} N. N. Vasiliev and D. A. Pavlov,
\emph{Enumeration of finite monomial orderings and combinatorics
of universal Gr\"obner bases}, Programming and Computer Software
35 (2009), 79--89.
\href{https://doi.org/10.1134/S0361768809020030}{doi:10.1134/S0361768809020030}.
\bibitem{M99} D. Maclagan,
\emph{Boolean term orders and the root system \(B_n\)}.
Inspected preprint version:
\href{https://arxiv.org/abs/math/9809134}{arXiv:math/9809134v2} (1999).
\bibitem{CCS} R. Christian, M. Conder and A. Slinko,
\emph{Flippable pairs and subset comparisons in comparative
probability orderings}. Inspected author-hosted manuscript,
\href{https://www.math.auckland.ac.nz/~slinko/Research/isipta17.pdf}{author copy}.
\bibitem{T14} T. Tvedebrink, P. S. Eriksen and N. Morling,
\emph{The multivariate Dirichlet-multinomial distribution and its
application in forensic genetics to adjust for sub-population
effects using the \(\theta\)-correction}.
\href{https://arxiv.org/abs/1406.6508}{arXiv:1406.6508v2} (2014),
Section 4.2 and Appendix A.
\bibitem{KB04} C. Kimberling and J. E. Brown,
\emph{Partial complements and transposable dispersions},
Journal of Integer Sequences 7 (2004), Article 04.1.6.
\href{https://cs.uwaterloo.ca/journals/JIS/VOL7/Kimberling/kimber67.pdf}{inspected article}.
\bibitem{K22} C. Kimberling,
\emph{Self-containing sequences, fractal sequences,
selection functions, and parasequences},
Journal of Integer Sequences 25 (2022), Article 22.2.1.
\href{https://cs.uwaterloo.ca/journals/JIS/VOL25/Kimberling/kimber16.pdf}{inspected article}.
\bibitem{V02} A. Vietri,
\emph{Combinatorial aspects of total weight orders over monomials
of fixed degree}, Order 19 (2002), 239--263.
\href{https://doi.org/10.1023/A:1021236722756}{doi:10.1023/A:1021236722756},
Proposition 6.2.
\bibitem{Mu03} I. M\"uller, \emph{Corner cuts and their polytopes},
Beitr\"age zur Algebra und Geometrie 44 (2003), 323--333.
\href{https://ftp.gwdg.de/pub/misc/EMIS/journals/BAG/vol.44/no.2/3.html}{inspected article}.
\bibitem{B08} Y. Berstein, H. Maruri-Aguilar, S. Onn,
E. Riccomagno and H. Wynn,
\emph{Minimal average degree aberration and the state polytope
for experimental designs}. Inspected preprint:
\href{https://arxiv.org/abs/0808.3055}{arXiv:0808.3055v1} (2008).
\bibitem{L21} A. Lemmens, \emph{Convexity of distinct sum sets}.
Inspected preprint:
\href{https://arxiv.org/abs/1812.09527}{arXiv:1812.09527v3} (2021).
\end{thebibliography}
\end{document}
